\section{Introduction}\label{sec:introduction}

A favorite site of a random walk is a site at which the accumulated local time
is maximal. Although the planar walk returns to the origin infinitely often,
its favorite sites need not remain near the origin. The question considered
here is how small a neighborhood of the origin can still contain a favorite
at arbitrarily large times. Because several sites can tie for the maximum,
there are two distinct observables: the distance of the nearest favorite and
the distance of the farthest favorite.

Let $(S_n)_{n\ge0}$ be simple symmetric random walk on $\mathbb Z^2$, started
at the origin. Define
\begin{equation}\label{eq:local-time-definition}
 L_n(x)=\sum_{j=0}^n\mathbf1_{\{S_j=x\}},\qquad
 M_n=\max_{x\in\mathbb Z^2}L_n(x),\qquad
 \mathcal F_n=\{x\in\mathbb Z^2:L_n(x)=M_n\}.
\end{equation}
The set $\mathcal F_n$ is finite and nonempty. Put
\begin{equation}\label{eq:radii-definition}
 R_n^- =\min_{x\in\mathcal F_n}\|x\|_2,\qquad
 R_n^+ =\max_{x\in\mathcal F_n}\|x\|_2,
 \qquad
 a_\gamma(n)=\frac{\sqrt n}{(\log n)^\gamma},\quad n\ge3.
\end{equation}
For example, along the path $(0,e_1,0,e_1)$ the two sites tie, so
$R_3^-=0$ and $R_3^+=1$. Thus localizing an arbitrarily chosen favorite
does not by itself localize the whole favorite set.

\begin{theorem}[Lower-class escape rate]\label{main:precise-escape-rate}
Almost surely, simultaneously for every real $\gamma$,
\begin{align}
 \liminf_{n\to\infty}\frac{R_n^-}{a_\gamma(n)}
 &=\infty, &&\gamma>\tfrac12,\label{eq:main-escape}\\
 \liminf_{n\to\infty}\frac{R_n^+}{a_\gamma(n)}
 &=0, &&\gamma\le\tfrac12.\label{eq:main-recurrence}
\end{align}
\end{theorem}

In particular, for each $\gamma>1/2$ and each fixed $A>0$, eventually
every favorite lies outside $B(0,Aa_\gamma(n))$. At the critical exponent,
for every $\varepsilon>0$ there are infinitely many times at which
\emph{all} favorites lie in $B(0,\varepsilon\sqrt{n/\log n})$.
The conclusion identifies the boundary among logarithmic powers of the
diffusive scale. It does not assert an integral test for arbitrary slowly
varying perturbations of the critical gauge.

Dembo, Peres, Rosen and Zeitouni~\cite[Theorem~1.1]{DPRZ2001} proved that
every planar favorite lies at distance $n^{1/2+o(1)}$. This first-order
exponent leaves the logarithmic scales in Theorem~\ref{main:precise-escape-rate}
undetermined. Erd\H{o}s and R\'ev\'esz
\cite[Question~14]{ErdosRevesz1987} asked whether planar favorite sites
escape to infinity. R\'ev\'esz later asked for their rate of escape
\cite[Chap.~19]{Revesz1990}, as recorded by Dembo~\cite[p.~6]{Dembo2005}.
Dembo's ICM survey~\cite[Open Problem~2.1]{Dembo2006} asks more broadly
about the evolution of favorite sites in dimensions $d\ge2$, and Hao,
Li, Okada and Zheng~\cite[Remark~1.4(1)]{HLOZ2025} also highlight the
planar escape-rate question. In one dimension the corresponding theorem
has critical exponent $1$; see Lifshits and Shi~\cite[Theorem~1.1]{LifshitsShi2004}
for the walk and Bass~\cite[Theorem~2.2]{Bass2023} for Brownian motion.
Our focus is the location of the full set of planar maximizers.
The two assertions in Theorem~\ref{main:precise-escape-rate} therefore retain
the distinction between its nearest and farthest members throughout the proof.

\subsection{Notation and the two local-time estimates}

All logarithms are natural. We write $P_x$ and $E_x$ for the law and
expectation of a walk started at $x$, and omit the subscript when $x=0$.
Its transition operator is
\[
 Pf(x)=\frac14\sum_{y:\|y-x\|_2=1}f(y).
\]
By symmetry the same notation acts on lattice measures. For $r>0$ let
$D(x,r)=\{y\in\mathbb Z^2:\|y-x\|_2<r\}$ and
$\tau_r=\inf\{n\ge0:S_n\notin D(0,r)\}$. Closed balls are denoted by
$B(x,r)$. We use the same geometric notation in the plane when the ambient
space is specified. A local time at a stopping time includes that time;
inside a killed domain it agrees with occupation strictly before exit.

The Green-function normalization and the local-time levels are
\begin{equation}\label{eq:level-definition}
 g=\frac2\pi,\qquad
 u_{J,c}=2gJ^2-cgJ\log J.
\end{equation}
The logarithmic radius is $J$, so the corresponding time scale is $e^{2J}$.
The coefficient of $J\log J$ will matter: knowing only the leading
$2gJ^2$ term cannot distinguish the powers in
Theorem~\ref{main:precise-escape-rate}.

The escape argument uses the following two estimates, proved in
Propositions~\ref{prop:localized-disk-upper} and
\ref{prop:disk-maximum-lower}. For fixed $\beta>0$, $c>2$ and $A>0$,
\begin{equation}\label{eq:intro-localized-upper}
 P\left(\max_{\|x\|_2\le e^JJ^{-\beta}}
 L_{\tau_{e^J}}(x)\ge u_{J,c}\right)
 \le J^{-2\beta+c-2+o(1)}+O(J^{-A}).
\end{equation}
For every fixed $\eta>0$, there is $p_\eta>0$ such that
\begin{equation}\label{eq:intro-disk-lower}
 P(M_{\tau_{e^J}}\ge u_{J,2+\eta})\ge p_\eta
 \quad\text{for all sufficiently large }J.
\end{equation}
Since $a_\gamma(e^{2J})=e^J/(2J)^\gamma$, the disk in
\eqref{eq:intro-localized-upper} is at the scale relevant to the theorem.
Taking $\beta=\gamma$ and $c=2+2\delta$ makes its main term
$J^{-2\gamma+2\delta+o(1)}$, summable over integer $J$ when
$\gamma>1/2$ and $\delta>0$ is small enough. The lower estimate is
amplified over independent segments to give an eventual almost-sure
bound at deterministic exponential times. Comparing it with the upper
estimate excludes a favorite near the origin throughout each time block.
Recurrence at the critical exponent uses a separate geometric-time
localization argument.

Constants denoted by $c,C$ may change between occurrences. Dependence on a
fixed parameter is indicated when relevant. Integer parts and fixed changes
of radius are harmless in estimates with a $J\log J$ margin: replacing $J$
by $J+O(1)$ changes $u_{J,c}$ by $O(J)$. Directed lattice crossings and
conditional bridge laws are defined exactly when they are introduced.
